% =====================================================================
% Appendix A — Versioning and Temporal Validity
% Task A7 (master plan Module 3.3)
% =====================================================================

\section{Versioning and Temporal Validity}
\label{app:versioning}

Large-language-model APIs are moving targets. A model ID served by a
vendor today may return different weights tomorrow, even though the
string in our code is unchanged. We document here the exact API strings
used, the invocation windows, and what we know about each model's
status during those windows. A machine-readable manifest mirroring this
section is archived at
\texttt{experiments\_log/models\_manifest.json}.

\subsection{Evaluated models (seven frontier LLMs)}

Table~\ref{tab:eval-models} lists the seven models whose outputs are
the object of evaluation in Experiments A and B.

\begin{table}[h]
\centering
\small
\renewcommand{\arraystretch}{1.2}
\caption{Evaluated models and invocation windows. ``Rolling alias''
indicates the vendor's model string is not a dated snapshot; weights
may change silently during the invocation window.}
\label{tab:eval-models}
\begin{tabular}{llll}
\toprule
\textbf{Paper name} & \textbf{API string} & \textbf{Vendor} & \textbf{Notes} \\
\midrule
Claude Opus 4.6     & \texttt{claude-opus-4-6}          & Anthropic    & Dated snapshot; immutable \\
Claude Sonnet 4.6   & \texttt{claude-sonnet-4-20250514} & Anthropic    & Dated snapshot; immutable \\
GPT-4o              & \texttt{gpt-4o}                   & OpenAI       & Rolling alias \\
OpenAI o3           & \texttt{o3}                       & OpenAI       & Rolling alias; reasoning family \\
Kimi v1-128k        & \texttt{moonshot-v1-128k}         & Moonshot AI  & Explicit version in ID \\
Gemini 2.5 Flash    & \texttt{gemini-2.5-flash}         & Google       & Rolling alias; uses ``thinking'' tokens \\
GPT-5               & \texttt{gpt-5}                    & OpenAI       & Rolling alias; reasoning family \\
\bottomrule
\end{tabular}
\end{table}

All seven models were invoked between March~2026 and April~2026, within
a single corpus (Experiment~A's 137 papers and Experiment~B's 100
synthetic scenarios).\footnote{GPT-5 completed 76 of 110 Experiment~A
papers before an API-quota interruption. After crediting, it resumed
and finished the remaining papers. We record the exact interruption
point in the run log at
\texttt{experiments\_log/runs/2026-04-22/001\_\_expa\_gpt5\_batch}.}
Rolling aliases (GPT-4o, o3, GPT-5, Gemini 2.5 Flash) introduce a known
limitation: a reader re-running our prompts against the same model ID
after a vendor update would observe different outputs. This is an
intrinsic constraint of API-only evaluation that no paper using these
endpoints can fully avoid. We mitigate by (a) recording invocation
windows, (b) saving every raw output at invocation time, and
(c) committing all model outputs into the release bundle.

\subsection{Annotator models (Phase 2 GT construction)}

Four LLMs cast votes on the canonical \texttt{method\_family} and
\texttt{conclusion\_direction} labels used for L1/L3/L4 scoring.
Table~\ref{tab:annot-models} summarizes.

\begin{table}[h]
\centering
\small
\renewcommand{\arraystretch}{1.2}
\caption{Phase 2 ground-truth annotator models. All four cast an equal
vote; the canonical label is the majority.}
\label{tab:annot-models}
\begin{tabular}{lll}
\toprule
\textbf{Role} & \textbf{API string} & \textbf{Vendor} \\
\midrule
Primary re-extractor (Phase~1 + Phase~2)  & \texttt{claude-opus-4-7} & Anthropic \\
Cached vote (original ingestion)          & \texttt{gpt-4o}          & OpenAI    \\
K2-generation voter                        & \texttt{kimi-latest}     & Moonshot AI \\
Independent voter                          & \texttt{gemini-2.5-flash}& Google    \\
\bottomrule
\end{tabular}
\end{table}

Three design choices deserve emphasis.

\paragraph{Primary re-extractor is outside the evaluated pool.}
Claude Opus~4.7 was chosen as the primary Phase~1 synthesis engine and
Phase~2 annotator precisely because the evaluated pool contains
\emph{Opus 4.6}, not 4.7. This avoids the degenerate case of a model
grading its own answers.

\paragraph{Kimi annotator was upgraded mid-voting.}
The initial Phase~2 vote used \texttt{moonshot-v1-128k}. We observed a
systematic failure mode (Event Study papers being mislabeled as Other)
and upgraded to \texttt{kimi-latest}, which auto-routes to Moonshot's
K2-generation model. The decision and its evidence are logged at
\texttt{experiments\_log/decisions/2026-04-22\_\_kimi\_upgrade.md}.
Only the final, upgraded votes are used for canonical labels.

\paragraph{GPT-4o appears in both pools.}
GPT-4o is both an evaluated subject (Experiments A and B) and an
annotator voter (Phase~2). This is a real circularity concern. Our
mitigations are (i)~GPT-4o is only one of four equally-weighted
annotator votes, (ii)~Module~2.2's Leave-One-Model-Out robustness
analysis explicitly recomputes rankings after removing each annotator
family, and (iii)~the primary re-extractor (Opus~4.7) is outside the
evaluated pool.

\subsection{Temporal scope}

All results in this paper should be read as characterizing these
specific models during the March--April 2026 window. We make no claim
about a given vendor's future releases. A reproducibility check against
the dated snapshot (where available) and a re-run against future rolling
aliases are both valid follow-ups. We commit to releasing raw outputs,
prompts, and the \texttt{models\_manifest.json} so that any divergence
can be attributed to a specific model change rather than a benchmark
ambiguity.
